\section{Scope, claim status, and the push-scale setting}
\label{sec:ladder-scope}

\subsection*{What this note does and does not claim}

This note fixes one instance, one state, one primitive, one work model, one
output representation, and one accuracy guarantee, and then proves how much
charged work is needed by three algorithm classes that differ by one axiom at
a time. It is a \emph{class} comparison. It is not an information lower
bound, and \Cref{sec:ladder-elimination} exhibits, on the same instance, an
algorithm outside the classes that is asymptotically faster than every member
of the strongest class by at least an $\Omega(1/\sqrt\alpha)$ factor. Nothing
in this note argues that $1/(\sqrt\alpha\,\epsppr)$ is a necessary cost of the
semantic problem.

\paragraph{Root problem and scope.}
The governing contract is Problem~1 and its graph-uniform open question in
\texttt{manuscript/notes/problem\_definitions/main.tex}. Thus
$\mathcal{G}$ is finite, simple, undirected, connected, unit-weighted, and has
$n\ge2$; the input is one seed vertex $v$ with
$\bm{s}=\bm{e}_v$ and $\alpha\in(0,1]$; and the requested object is a sparse
approximation to $\bm{x}_0^*$. This note studies the star subfamily and
$0<\epsppr\le1/16$, chiefly to test whether restricted one-hop algorithm
classes can attain the open fully charged target
$\widetilde{O}(1/(\sqrt\alpha\,\epsppr))$. It neither studies RPPR Problem~2
nor the nested-SDD reuse question, and it proves no graph-uniform lower bound
for the semantic PPR problem.

\paragraph{Accuracy namespace.}
Throughout, $\epsppr$ is exactly the semantic degree-normalized PPR output
accuracy of Problem~1. The source-native \APPR{} activation threshold
$\epsappr$ appears only in \Cref{sec:ladder-manuscript-relation}; assigning it
the same numerical value as $\epsppr$ there is an explicit comparison choice,
not an identification of the two meanings. No RPPR objective or KKT tolerance
is used in this note.

\paragraph{Claim status.}
Statements carry one of the following labels.
\statusproveddraft{} means a complete argument is written out in this note but
has not been independently audited; it asserts no repository promotion.
\statusimported{} means the argument is complete in a named source package and
is used here without being re-derived. \statusmeasured{} means numerics only.
\statusopen{} means open. Provenance for every statement, together with an
auditor's checklist, is in \Cref{sec:ladder-verification}.

\subsection*{Push-scale notation}

The algebraic definitions in
\Cref{subsec:shared-source-aligned-problem} are used verbatim. The current
Problem~1 output and fully charged work contract, stated above, governs every
complexity claim in this note. This subsection only introduces note-scoped
abbreviations and a mass-coordinate state; it redefines no shared object. Put
\begin{equation}
    \calpha:=\frac{1-\alpha}{1+\alpha},
    \qquad
    \galpha:=\frac{2\alpha}{1+\alpha},
    \qquad
    \bm{H}:=\bm{I}-\calpha\bm{A}\bm{D}^{-1},
    \label{eq:ladder-push-scale-constants}
\end{equation}
so that $1-\calpha=\galpha$ and $1+\calpha=2/(1+\alpha)$. Conjugating the
shared operator $\bm{Q}$ of \eqref{eq:shared-pagerank-matrices} gives
\begin{equation}
    \bm{D}^{1/2}\bm{Q}\bm{D}^{-1/2}
    =\frac{1+\alpha}{2}\,\bm{H},
    \qquad\text{hence}\qquad
    \bm{H}\bm{\pi}=\galpha\bm{e}_{v}
    \label{eq:ladder-H-conjugation}
\end{equation}
for the single-seed specialization $\bm{s}=\bm{e}_{v}$ of
\eqref{eq:shared-seed}. Write
\begin{equation}
    \pr(\bm{h}):=\galpha\bm{H}^{-1}\bm{h},
    \qquad\text{so}\qquad
    \pr(\bm{e}_{v})=\bm{\pi} .
    \label{eq:ladder-pr-operator}
\end{equation}
The symmetric form used in \Cref{sec:ladder-signed-lower-bound} is
\begin{equation}
    \Hsym
    :=\bm{I}-\calpha\bm{D}^{-1/2}\bm{A}\bm{D}^{-1/2}
    =\frac{2}{1+\alpha}\,\bm{Q},
    \qquad
    \galpha\bm{I}\preceq\Hsym\preceq\frac{2}{1+\alpha}\bm{I},
    \label{eq:ladder-Hsym}
\end{equation}
the spectral bracket following from
$\alpha\bm{I}\preceq\bm{Q}\preceq\bm{I}$.

\subsection*{State, primitive, and fully charged work}

An algorithm maintains a pair $(\bm{z},\bm{r})$ with
\begin{equation}
    \bm{z}^{0}=\bm{0},
    \qquad
    \bm{r}:=\galpha\bm{e}_{v}-\bm{H}\bm{z},
    \qquad\text{so}\qquad
    \bm{r}^{0}=\galpha\bm{e}_{v},
    \qquad
    \bm{\xi}:=\bm{\pi}-\bm{z}=\bm{H}^{-1}\bm{r} .
    \label{eq:ladder-state}
\end{equation}
For components, write $z_u:=(\bm{z})_u$, $r_u:=(\bm{r})_u$,
$\xi_u:=(\bm{\xi})_u$, and $\pi_u:=(\bm{\pi})_u$. The one-hop primitive is, for
$u\in\mathcal{V}$ and $\delta\in\R$,
\begin{equation}
    \push(u,\delta):\quad
    z_{u}\mathrel{+}=\delta,
    \qquad
    r_{u}\mathrel{-}=\delta,
    \qquad
    r_{w}\mathrel{+}=\frac{\calpha\,\delta}{d_{u}}
    \quad(w\in\mathcal N(u)),
    \label{eq:ladder-primitive}
\end{equation}
and it inspects all $d_u$ entries of $u$'s adjacency list. By construction,
$\push(u,\delta)$ preserves
$\bm{z}+\bm{H}^{-1}\bm{r}=\bm{\pi}$ exactly and changes $\bm{z}$ only at its
base vertex; the neighbor split in \eqref{eq:ladder-primitive} is forced and
the scalar $\delta$ is the only freedom.

Write $W_{\mathrm{adj}}:=\sum_t d_{u_t}$ for adjacency-entry inspections and
let $\Work$ denote the \emph{fully charged} exact-real word work of Problem~1:
it also counts arithmetic, comparisons, state reads and writes, degree
queries, and emitted output words. Thus each push costs $\Theta(d_u)$ in the
full model and $\Work\ge W_{\mathrm{adj}}$. All lower bounds below first
lower-bound $W_{\mathrm{adj}}$ and therefore also lower-bound $\Work$; stated
upper bounds perform $O(1)$ additional word operations per inspected entry or
emitted coordinate. Preprocessing is not free. The policy choosing
$(u_t,\delta_t)$ is arbitrary --- adaptive, randomized, or batched. Every
lower bound is pathwise, so it holds for every realized random path; a claimed
randomized upper bound would additionally have to state success probability
and the type of work guarantee.

\subsection*{Semantic output guarantee}

The algorithm reports a sparse list for $\widehat{\bm{x}}$, with unlisted
coordinates equal to zero, exactly as in Problem~1. In the mass coordinates
used by this note, write
\begin{equation}
    \widehat{\bm{\pi}}:=\widehat{\bm{z}},
    \qquad
    \widehat{\bm{x}}:=\bm{D}^{-1/2}\widehat{\bm{z}},
    \qquad
    \errop(\widehat{\bm{z}})
    :=\max_{u\in\mathcal{V}}
      \frac{\abs{\pi_{u}-\widehat z_{u}}}{d_{u}}
    =\norm{\bm{D}^{-1/2}
      (\widehat{\bm{x}}-\bm{x}_0^*)}_{\infty}
    \le\epsppr .
    \label{eq:ladder-guarantee}
\end{equation}
The two lists have the same support; converting a listed mass coordinate to
an $\bm{x}$-coordinate and emitting it are charged operations.

\subsection*{Dictionary to the source mass scale}

The monotone results were first proved in the mass scale of
\file{sections/appr_lower_bound.tex}, where the state is the \APPR{} pair
$(\bm{p},\bm{r}_{\mathrm{mass}})$. The translation is
\begin{equation}
    \bm{z}=\bm{p},
    \qquad
    \bm{r}=\galpha\bm{r}_{\mathrm{mass}},
    \qquad
    \delta=\alpha\eta,
    \qquad
    \bm{\pi}-\bm{z}=\bm{H}^{-1}\bm{r}
    =\pr(\bm{r}_{\mathrm{mass}}).
    \label{eq:ladder-mass-dictionary}
\end{equation}
Under it, lazy \APPR{} push at $u$ is
$\push(u,\omega r_u)$ with $\omega=(1+\alpha)/2$; full non-lazy push is
$\omega=1$; and Gauss--Seidel SOR with $\omega>1$ overshoots $r_u$ past
zero. Feasibility of the monotone class,
$\eta\le2(\bm{r}_{\mathrm{mass}})_u/(1+\alpha)$, is exactly
$\delta\le r_u$.

\subsection*{Five facts used throughout}

Let $\bm{d}:=(d_u)_{u\in\mathcal{V}}$. The following elementary identities
keep the push scale and the source mass scale separate.

\begin{lemma}[Elementary facts; \statusproveddraft]
\label{lem:ladder-facts}
For every connected $\mathcal{G}$, every $\alpha\in(0,1)$, and every
$\bm{h}\in\R^{n}$:
\begin{enumerate}[label=\textup{(F\arabic*)},leftmargin=3em]
    \item\label{it:ladder-f1}
        $\bm{H}^{\trans}\one=\galpha\one$, hence
        $\sum_u(\bm{H}^{-1}\bm{r})_u
        =\norm{\bm{r}}_1/\galpha$ whenever $\bm{r}\ge\bm{0}$;
    \item\label{it:ladder-f2}
        $\bm{H}^{-1}\ge\bm{0}$ entrywise;
    \item\label{it:ladder-f3}
        $(\bm{H}^{-1})_{vv}=\pi_v/\galpha$;
    \item\label{it:ladder-f4}
        $\max_u\pr(\bm{h})_u/d_u\le\max_u h_u/d_u$ for
        $\bm{h}\ge\bm{0}$;
    \item\label{it:ladder-f5}
        $\abs{\bm{H}^{-1}\bm{h}}
        \le\bm{H}^{-1}\abs{\bm{h}}$ entrywise.
\end{enumerate}
\end{lemma}

\begin{proof}
\ref{it:ladder-f1}:
$\bm{H}^{\trans}\one
=\one-\calpha\bm{D}^{-1}\bm{A}\one
=\one-\calpha\bm{D}^{-1}\bm{d}
=(1-\calpha)\one=\galpha\one$. Therefore
$\sum_u(\bm{H}^{-1}\bm{r})_u
=\inner{\bm{H}^{-\trans}\one}{\bm{r}}
=\norm{\bm{r}}_1/\galpha$ when $\bm{r}\ge\bm{0}$.
\ref{it:ladder-f2}: $\calpha\bm{A}\bm{D}^{-1}\ge\bm{0}$ has column sums
$\calpha<1$, so its Neumann series converges to $\bm{H}^{-1}$ and is
nonnegative.
\ref{it:ladder-f3}: take the $v$-coordinate of
$\bm{\pi}=\galpha\bm{H}^{-1}\bm{e}_v$.
\ref{it:ladder-f4}:
$\bm{H}\bm{d}=\bm{d}-\calpha\bm{A}\one=\galpha\bm{d}$, so
$\pr(\bm{d})=\bm{d}$; if $\bm{h}\le t\bm{d}$ coordinatewise then
$\pr(\bm{h})\le t\bm{d}$ by \ref{it:ladder-f2}.
\ref{it:ladder-f5}: immediate from \ref{it:ladder-f2}.
\end{proof}
